\documentclass{article}
\usepackage{amsmath} 
\usepackage{amsfonts}
\usepackage{amssymb}
\usepackage[parfill]{parskip}
\usepackage{esvect}
\usepackage[thinc]{esdiff}
\usepackage{mathtools}
\usepackage{physics}
\usepackage{amsthm}
\usepackage{xfrac}
\usepackage{enumitem}

\DeclareMathOperator{\im}{im}

\begin{document}

\title{MATH 430: HW 5}
\author{Jacob Lockard}
\date{5 March 2026}

\maketitle

\begin{quote}
    \itshape
    For the exercises below, assume $(G, *_G)$ and $(H, *_H)$ are groups,
    and $\phi : G \to H$ is a homomorphism. Write $e_G$ and $e_H$
    for the identity elements in $G$ and $H$, respectively.
\end{quote}

We omit the operators $*_G$ and $*_H$ throughout this text, using the typical
multiplication and exponent notation for groups.

\section*{Exercise 1}

We'll show that $\phi(e_G) = e_H$. We have:
\begin{align*}
    \phi(e_G) = \phi(e_G e_G) = \phi(e_G) \phi(e_G) = \phi(e_G)^2 \,.
\end{align*}
We've proven previously that a group's only idempotent is its identity.
So $\phi(e_G)$ is the identity in $H$.
\qedsymbol

\section*{Lemma 1}

Let $x \in G$. We show that $\phi(x)^{-1} = \phi(x^{-1})$:
\begin{align*}
    \phi(x)\phi(x^{-1}) = \phi(xx^{-1}) = \phi(e_G) = e_H =
    \phi(e_G) = \phi(x^{-1}x) = \phi(x^{-1})\phi(x) \,,
\end{align*}
by the previous exercise and uniqueness of group inverses.
\qedsymbol

\section*{Exercise 2}

We'll show that $\im(\phi) \leq H$. By definition, $\im(\phi) \subseteq H$.
Exercise 1 ensures $e_H \in \im(\phi)$.
Let $\phi(a), \phi(b) \in \im(\phi)$.
Then by the previous lemma,
\begin{align*}
    \phi(a) \phi(b)^{-1} = \phi(a) \phi(b^{-1}) = \phi(ab^{-1}) \in \im(\phi) \,.
\end{align*}
So $\im(\phi) \leq H$.
\qedsymbol

\section*{Exercise 3}

Suppose $G$ is abelian. We'll show that $\im(\phi)$ is also abelian.
Let $\phi(a), \phi(b) \in \im(\phi)$. Then,
\begin{align*}
    \phi(a)\phi(b) = \phi(ab) = \phi(ba) = \phi(b)\phi(a) \,.
\end{align*}
\qedsymbol

On the other hand, we note that $H$ need not be abelian. For example, let $G = \{\iota\} \subseteq D_3$
and $H = D_3$, where $\iota$ is the identity in $D_3$; and let $\phi$
map $\iota$ to $\iota$. Then $\phi$ is clearly a homomorphism and clearly $\iota\iota = \iota\iota$,
but we know $D_3$ is not abelian.

\section*{Exercise 4}

Suppose $\phi$ is an isomorphism.
We'll show that $G$ is abelian if and only if
$H$ is abelian.

If $G$ is abelian,
then the previous exercise ensures
$\im(\phi)$ is abelian.
But since $\phi$ is surjective, its image is $H$, so $H$ is abelian.

If $G$ is not abelian, then there's an $a, b \in G$ such that
$ab \neq ba$. Since $\phi$ is injective, we have
$\phi(a)\phi(b) = \phi(ab) \neq \phi(ba) = \phi(b)\phi(a)$.
So $H$ is not abelian.
\qedsymbol

\section*{Exercise 5}

We'll show that $\ker(\phi) \leq G$. By definition, $\ker(\phi) \subseteq G$.
Exercise 1 ensures $e_G \in \ker(\phi)$.
Let $a, b \in \ker(\phi)$.
Then by Lemma 1,
\begin{align*}
    \phi(ab^{-1}) = \phi(a) \phi(b^{-1}) = \phi(a) \phi(b)^{-1} = e_H e_H^{-1} = e_H \,,
\end{align*}
So $ab^{-1} \in \ker(\phi)$. We conclude that $\ker(\phi) \leq G$.
\qedsymbol

\section*{Exercise 6}

We'll show that $\ker(\phi)$
is a normal subgroup of $G$. We've already shown $\ker(\phi) \leq G$.
Let $g \in G$ and $n \in \ker(\phi)$. Then,
\begin{align*}
    \phi(gng^{-1}) &= \phi(g) \phi(n) \phi(g^{-1}) \\
    &= \phi(g) \phi(g^{-1}) \\
    &= \phi(g g^{-1}) \\
    &= \phi(e_G) = e_H \,,
\end{align*}
where the last line is justified by Exercise 1.
So $gng^{-1} \in \ker(\phi)$.
We conclude that $\ker(\phi)$ is a normal subgroup of $G$. \qedsymbol

\section*{Exercise 7}

We leave it to the reader to show that
that $G \times H$ is a group under the operation
defined by
\begin{align*}
    (g,h) (g', h') := (gg', hh') \,,
\end{align*}
for all $(g, h), (g', h') \in G \times H$. We simply note that
the identity in $G \times H$ is $(e_G, e_H)$, and that for any $(g, h) \in G \times H$,
\begin{align*}
    (g, h)^{-1} = (g^{-1}, h^{-1}) \,.
\end{align*}

\section*{Exercise 8}

Let $\pi_G : G \times H \to G$
be the projection $(g,h) \mapsto g$.
We'll show that $\pi_G$ is a group
homomorphism. Let $(g, h), (g', h') \in G \times H$. We have:
\begin{align*}
    \pi_G \bigl( (g, h) (g', h') \bigr) = \pi_G (gg', hh') = gg'
    = \pi_G(g, h) \pi_G(g', h') \,.
\end{align*}

\section*{Exercise 9}

We'll show that $G \times \{e_H\}$
is a normal subgroup of $G \times H$.
Clearly, $G \times \{e_H\} \subseteq G \times H$ and
$(e_G, e_H) \in G \times \{e_H\}$.
Let $(g, e_H), (g', e_H) \in G \times \{e_H\}$. Then,
\begin{align*}
    (g, e_H) (g', e_H)^{-1} &= (g, e_H) ((g')^{-1}, e_H) \\
    &= (g(g')^{-1}, e_H e_H) \in G \times \{e_H\} \,.
\end{align*}
So $G \times \{e_H\} \leq G \times H$.
Now let $(g, e_H) \in G \times \{e_H\}$ and $(g', h') \in G \times H$.
Then,
\begin{align*}
    (g', h') (g, e_H) (g', h')^{-1} &= (g', h')(g, e_H)((g')^{-1}, (h')^{-1}) \\
    &= \bigl(g'g(g')^{-1}, h'e_H(h')^{-1}\bigr) \\
    &= \bigl(g'g(g')^{-1}, e_H\bigr) \in G \times \{e_H\} \,.
\end{align*}
So $G \times \{e_H\}$ is a normal subgroup of $G \times H$.
\qedsymbol

\section*{Exercise 10}

Suppose $g \in G$ has finite order.
We'll show $\phi(g)$ has finite order. We have, for some $n \in \mathbb{Z}^+$,
\begin{align*}
    g^n  &= e_G \\
    \phi(g)^n = \phi(g^n) &= \phi(e_G) = e_H \,.
\end{align*}
\qedsymbol

\section*{Challenge}

Let $g \in G$ have order $n$.
We'll show that $m|n$, where $m$ is the order of $\phi(g)$.
We first note that
\begin{align*}
    \phi(g)^n = \phi(g^n) = \phi(e_G) = e_H \,,
\end{align*}
so $m \leq n$, since $m$ is the smallest positive integer such that $\phi(g)^m = e_H$.
The discrete division theorem ensures we can write
\begin{align*}
    n = qm + r & \,,
\end{align*}
for some unique $q, r \in \mathbb{Z}$ with $0 \leq r < m$.
If $q \leq 0$, then since $m > 0$,
\begin{align*}
    qm &\leq 0 \\
    n = qm + r &\leq r < m \leq n \,.
\end{align*}
Since $n = n$, we conclude that $q > 0$.
We have:
\begin{gather*}
    g^n = g^{qm + r} = g^{qm} g^r \\
    g^{n-qm} = g^r \,.
\end{gather*}
Since $q$ and $m$ are positive, we know $n-qm < n$, so
$g^{n-qm} \neq e_G$, since $n$ is the smallest positive integer
such that $g^n = e_G$.
We conclude $r \neq 0$.
So we can write $n = qm$, and we've shown that $m|n$.
\qedsymbol

\end{document}